% !TeX root = ../../Book2.tex
\chapter{Hermite、Legendre 与球谐函数}
\label{b2:app:B}

\begin{chapterabstract}
本附录构造三组常用特征函数。Hermite 多项式乘 Gauss 权形成实线上的正交基，同时对角化谐振子与适当归一化的 Fourier 变换。Legendre 多项式解决区间上的奇异端点方程，并对应三维球面的轴对称模态。一般球谐函数则由齐次调和多项式构造；通过调和分解和多项式逼近证明其完备性，进一步得到加法公式、球内调和延拓及球形 Helmholtz 方程的径向 Bessel 阶数。
\end{chapterabstract}
\begin{learninggoals}
\begin{itemize}
\item\label{b2:item:B2-goal-hermite}从 Rodrigues 公式证明 Hermite 正交性与完备性，并核对两种常用 Gauss 归一化。
\item\label{b2:item:B2-goal-oscillator}用升降算子求谐振子谱，区分微分表达式与自伴算子定义域。
\item\label{b2:item:B2-goal-legendre}由正交结构推导 Legendre 方程、范数和递推，保留端点及收敛意义。
\item\label{b2:item:B2-goal-spherical}构造球谐空间并证明谱完备，联系极坐标 Laplace 算子与 Bessel 径向方程。
\end{itemize}
\end{learninggoals}

分离变量法的第一步是找到若干特征函数，最后一步却是说明这些函数是否足以表示一般数据。二者之间有不可省略的完备性证明。本附录选取三种可以把这一证明完整写出的模型：实线上的 Gauss 权、有限区间上的无权积分，以及紧球面上的面积测度。它们各自的几何与权重决定了归一化和特征值。

\ProseParagraphBreak

正文只用本卷的 Fourier 唯一性、Plancherel 与常微分方程基础，并复用附录 A 的 Gamma、Beta 积分。相关特殊函数的标准公式可参阅 \cite[Chapters 14 and 18]{NationalInstituteofStandardsNIST}；这里把正交、完备、算子作用与应用逐项展开，不将表中公式代替函数展开所需的证明。

% !TeX root = ../../Book2.tex
\section{Hermite 函数与谐振子}
\label{b2:sec:B201}

指数函数适合平移不变算子，带有二次势能的算子却同时涉及位置和频率。Hermite 函数把多项式与 Gauss 衰减结合起来，既形成全空间的正交基，又使谐振子和 Fourier 变换具有简单的谱表示。归一化必须与 Fourier 指数中的 \(2\pi\) 分开核对。

\subsection{多项式、正交性与完备性}
\label{b2:subsec:B20101}

定义 \term[hermite-duoxiangshi]{Hermite 多项式}[Hermite polynomial]
\[H_n(x)=(-1)^ne^{x^2}\frac{\mathrm d^n}{\mathrm dx^n}e^{-x^2},\qquad n=0,1,2,\ldots.\]
这是权重为 \(e^{-x^2}\) 的常用约定；概率论中权重 \(e^{-x^2/2}\) 的同名多项式需要缩放。由定义求导得到 \(H_{n+1}=2xH_n-H_n'\)，故 \(H_n\) 是首项 \(2^nx^n\) 的 \(n\) 次多项式。前几项为 \(1,2x,4x^2-2,8x^3-12x\)。

\begin{proposition}[Hermite 正交基]
\label{b2:prop:B2-hermite-basis}
令
\[\widetilde h_n(x)=\frac{H_n(x)e^{-x^2/2}}{(2^nn!\sqrt\pi)^{1/2}}.\]
则 \((\widetilde h_n)_{n\ge0}\) 是 \(L^2(\mathbb R)\) 的标准正交基。
\end{proposition}
\begin{proof}
分部积分 \(n\) 次，所有边界项都因多项式乘 Gauss 衰减而为零，得
\[\int_\mathbb R H_nH_m e^{-x^2}\,\mathrm dx
 =\int_\mathbb R e^{-x^2}H_m^{(n)}\,\mathrm dx.
\]
若 \(m<n\)，右端为零；\(m=n\) 时为 \(2^nn!\sqrt\pi\)。交换指标处理另一半，得到标准正交性。

为证完备性，设 \(g\in L^2\) 与全部 \(\widetilde h_n\) 正交。由于 \(H_n\) 的阶数逐次增加且首系数非零，它们张成全部多项式，故 \(\int g(x)e^{-x^2/2}x^n\,\mathrm dx=0\) 对每个 \(n\) 成立。对固定实数 \(t\)，Cauchy--Schwarz 不等式保证
\[\int|g(x)|e^{-x^2/2}e^{|t||x|}\,\mathrm dx<\infty.\]
因此可以把 \(e^{-itx}\) 的幂级数逐项积分，得到 \(g e^{-x^2/2}\) 的 Fourier 变换在所有实频率为零。该乘积属于 \(L^1\)，本卷已经证明的 Fourier 唯一性推出它几乎处处为零，继而 \(g=0\)。正交补为零即完备。
\end{proof}

函数 \(\widetilde h_n\) 是另一种尺度归一化的 Hermite 函数；与第五章的归一化之间的关系将在本节末写明。证明中的正交性来自反复分部积分，完备性却依赖对全部频率的唯一性；前者不能替代后者。仅计算出许多正交模态，尚不足以断言它们可以展开任意平方可积函数。
\symidx{\(\widetilde h_n\)：谐振子尺度的 Hermite 函数}

\subsection{升降算子与二次势能}
\label{b2:subsec:B20102}

在 Schwartz 函数上定义
\[a=\frac1{\sqrt2}\left(x+\frac{\mathrm d}{\mathrm dx}\right),\qquad
 a^*=\frac1{\sqrt2}\left(x-\frac{\mathrm d}{\mathrm dx}\right).
\]
星号与 \(L^2\) 伴随关系相容，因为分部积分把微分变号。直接展开给出 \(aa^*-a^*a=I\)，以及
\[\mathcal L=-\frac{\mathrm d^2}{\mathrm dx^2}+x^2=2a^*a+I.\]
算子 \(\mathcal L\) 称为 \term[xiezhenzi]{谐振子}[harmonic oscillator]；此处固定无量纲系数，物理质量与频率可通过缩放恢复。

\begin{proposition}[谐振子的特征函数]
\label{b2:prop:B2-oscillator-eigenfunctions}
Hermite 函数满足
\[a\widetilde h_n=\sqrt n\,\widetilde h_{n-1},\qquad a^*\widetilde h_n=\sqrt{n+1}\,\widetilde h_{n+1},\qquad
 \mathcal L\widetilde h_n=(2n+1)\widetilde h_n,\]
其中 \(a\widetilde h_0=0\)。
\end{proposition}
\begin{proof}
由 \(H_{n+1}=2xH_n-H_n'\) 直接得到升算子公式。又 \(a\widetilde h_0=0\)，而 \([a,a^*]=I\) 推出 \(a(a^*)^n=(a^*)^na+n(a^*)^{n-1}\)。反复使用升算子公式得到 \(\widetilde h_n=(a^*)^n\widetilde h_0/\sqrt{n!}\)，代入即得降算子公式。最后由 \(\mathcal L=2a^*a+I\) 得到特征值。
\end{proof}

对任意 \(f\in L^2\)，写 \(f=\sum c_n\widetilde h_n\)。可以精确定义自伴实现为
\[
 D(\mathcal L)=\left\{f:\sum_{n=0}^\infty(2n+1)^2|c_n|^2<\infty\right\},\qquad
 \mathcal Lf=\sum_{n=0}^\infty(2n+1)c_n\widetilde h_n.
\]
这个定义域不是全部 \(L^2\)。系数表示把算子变成实数序列的乘法；伴随泛函有界的充要条件恰是同一个加权平方可和条件，因此它自伴。对 Schwartz 函数，分部积分与特征函数公式说明该实现等于微分表达式。有限 Hermite 和在图范数中稠密，因此由算子在有限和上的作用及其闭性，就能恢复整个自伴实现。

\ProseParagraphBreak

逆算子 \((\mathcal L+I)^{-1}\) 的特征值为 \((2n+2)^{-1}\to0\)，截断为前 \(N\) 项后，算子范数误差为 \((2N+4)^{-1}\)。因此该逆是紧算子。二次势能把函数的质量限制在有限位置，离散特征值趋向无穷；这与全空间自由 Laplace 算子的连续 Fourier 频率形成不同的谱结构。

\subsection{本卷 Fourier 约定下的特征基}
\label{b2:subsec:B20103}

对本卷的 \(2\pi\) Fourier 变换，定义缩放后的标准正交基
\[\psi_n(x)=(2\pi)^{1/4}\widetilde h_n(\sqrt{2\pi}\,x).\]
其中 \(\psi_0(x)=2^{1/4}e^{-\pi x^2}\)，正好由 Gauss Fourier 公式保持不变。这里的 \(\psi_n\) 就是第五章记作 \(h_n\) 的归一化；本附录用 \(\widetilde h_n\) 区分谐振子尺度，避免在两种高斯指数之间混用公式。

\begin{theorem}[Fourier 变换的 Hermite 特征展开]
\label{b2:thm:B2-fourier-hermite}
对所有 \(n\ge0\)，有 \(\widehat\psi_n=(-i)^n\psi_n\)。因此若 \(f=\sum c_n\psi_n\) 于 \(L^2(\mathbb R)\)，则
\[\widehat f=\sum_{n=0}^\infty(-i)^nc_n\psi_n\quad\text{于 }L^2.\]
\end{theorem}
\begin{proof}
缩放后的升算子与 \(2\pi x-\frac{\mathrm d}{\mathrm dx}\) 成正比。Fourier 微分公式给出
\[\mathcal F\left(\left(2\pi x-\frac{\mathrm d}{\mathrm dx}\right)f\right)
 =-i\left(2\pi\xi-\frac{\mathrm d}{\mathrm d\xi}\right)\widehat f.
\]
从 \(\widehat\psi_0=\psi_0\) 出发，逐次使用升算子得到全部特征值。对有限和结论已成立，利用 Fourier 的 \(L^2\) 连续性和正交展开取极限，即得一般形式。
\end{proof}

由此再次看到 \(\mathcal F^4=I\)，以及偶、奇模态分别对应反射的正、负特征值。但每个 Fourier 特征值都有无穷多个模态：\(n\) 相差 \(4\) 时相位相同，故 Fourier 变换的四个特征空间都无限维。谐振子的特征值则逐阶区分这些模态，两个谱展开虽共用一组函数，却表达不同算子的信息。

\ProseParagraphBreak

在 \(\mathbb R^d\) 上取张量积 \(\widetilde h_\alpha(x)=\prod_{j=1}^d \widetilde h_{\alpha_j}(x_j)\)。Fubini 和一维完备性逐坐标应用，给出 \(L^2(\mathbb R^d)\) 的标准正交基。算子 \(-\Delta+|x|^2\) 的特征值为 \(2|\alpha|+d\)，所以第 \(n\) 层的重数是满足 \(|\alpha|=n\) 的多重指标个数 \(\binom{n+d-1}{d-1}\)。这一简并来自不同坐标方向分配同一总阶数，与球谐展开按角动量分类的方式不同。

% !TeX root = ../../Book2.tex
\section{Legendre 多项式与角向方程}
\label{b2:sec:B202}

球面上的轴对称函数只依赖极角，令 \(t=\cos\theta\) 后，其微分方程变成带系数 \(1-t^2\) 的一维方程。这个系数在两端为零，反映了球坐标在两极的退化。Legendre 多项式既给出规则解，也形成区间上无权测度的正交基。

\subsection{Rodrigues 公式与范数}
\label{b2:subsec:B20201}

定义 \term[legendre-duoxiangshi]{Legendre 多项式}[Legendre polynomial]
\[P_n(t)=\frac1{2^nn!}\frac{\mathrm d^n}{\mathrm dt^n}(t^2-1)^n.\]
它的首系数为 \((2n)!/[2^n(n!)^2]\)，并满足 \(P_n(1)=1\)、\(P_n(-t)=(-1)^nP_n(t)\)。第一个端点值可直接在 \((t-1)^n(t+1)^n\) 的 \(n\) 阶导数中取 \(t=1\)：只有把全部 \(n\) 次微分作用于 \((t-1)^n\) 的项保留下来。

\begin{proposition}[Legendre 正交基]
\label{b2:prop:B2-legendre-orthogonal}
有
\[\int_{-1}^1P_n(t)P_m(t)\,\mathrm dt=\frac2{2n+1}\delta_{nm}.\]
归一化函数 \(\sqrt{(2n+1)/2}\,P_n\) 构成 \(L^2(-1,1)\) 的标准正交基。
\end{proposition}
\begin{proof}
对次数小于 \(n\) 的多项式 \(q\)，把 Rodrigues 公式分部积分 \(n\) 次。\((t^2-1)^n\) 在两端具有 \(n\) 阶零点，所以全部边界项为零；最后 \(q^{(n)}=0\)，得到 \(\int P_nq=0\)。这给出不同阶数正交。若 \(q=P_n\)，则 \(P_n^{(n)}=(2n)!/(2^nn!)\)，所以
\[\int_{-1}^1P_n^2=\frac{(2n)!}{2^{2n}(n!)^2}\int_{-1}^1(1-t^2)^n\,\mathrm dt.
\]
附录 A 的 Beta 积分给出后一积分为 \(2^{2n+1}(n!)^2/(2n+1)!\)，于是范数等于 \(2/(2n+1)\)。

由于阶数和首系数不退化，\(P_n\) 张成全部多项式。多项式在 \(C[-1,1]\) 中一致稠密，例如把区间线性变到 \([0,1]\)，使用 Bernstein 多项式；其权重和为 \(1\)，离均值超过 \(\varepsilon\) 的权重总量由方差 \(t(1-t)/N\) 控制，因此一致连续性给出一致逼近。连续函数又在 \(L^2\) 中稠密，故正交系完备。
\end{proof}

前几项为 \(P_0=1\)、\(P_1=t\)、\(P_2=(3t^2-1)/2\)、\(P_3=(5t^3-3t)/2\)。与 Hermite 情形相比，这里区间有限且权重为 \(1\)，归一化常数完全不同。同样叫作“正交多项式”，并不意味着可以交换其积分区间与权重。
\symidx{\(P_n(t)\)：满足 \(P_n(1)=1\) 的 Legendre 多项式}

\subsection{微分方程与三项递推}
\label{b2:subsec:B20202}

\begin{proposition}[Legendre 方程]
\label{b2:prop:B2-legendre-equation}
Legendre 多项式满足
\[(1-t^2)P_n''-2tP_n'+n(n+1)P_n=0,\]
以及对 \(n\ge1\) 的递推
\[(2n+1)tP_n=(n+1)P_{n+1}+nP_{n-1}.
\]
\end{proposition}
\begin{proof}
令 \(Lq=-((1-t^2)q')'\)。它把次数不超过 \(n\) 的多项式映到同一空间，并把首项 \(a_nt^n\) 映成首项 \(n(n+1)a_nt^n\)。故 \(Q=LP_n-n(n+1)P_n\) 次数小于 \(n\)。对任意低次多项式 \(q\)，两次分部积分时边界项都由 \(1-t^2=0\) 消失，得 \(\int Qq=\int P_nLq-n(n+1)\int P_nq=0\)。取 \(q=Q\) 得 \(Q=0\)，即微分方程。

展开 \(tP_n\) 于 \(P_0,\ldots,P_{n+1}\)。对 \(m\le n-2\)，\(\int tP_nP_m=\int P_n(tP_m)=0\)，因为 \(tP_m\) 次数小于 \(n\)；奇偶性又排除 \(P_n\) 项。因此只剩 \(P_{n+1}\) 与 \(P_{n-1}\)。比较首系数得前者系数 \((n+1)/(2n+1)\)，在 \(t=1\) 取值使两系数和为 \(1\)，得到后者系数 \(n/(2n+1)\)。
\end{proof}

这个证明没有对 Rodrigues 公式作冗长展开，而是利用对称微分算子保持多项式次数的事实。正交性、微分方程与递推不是三张互不相关的公式表；它们来自同一有限维多项式结构。

\ProseParagraphBreak

方程的端点仍须认真处理。\(n=0\) 时除常数解外，还有 \(\frac12\log((1+t)/(1-t))\)，它在 \(t=\pm1\) 发散。两极对应球面上普通点，平滑角向函数不能容纳这种坐标奇性。因此在球面应用中选择规则的多项式解，是几何正则性而非任意约定。若改在不含端点的子区间求初值问题，第二支完全合法，不能因为它不是多项式就将其删除。

\subsection{Legendre 展开及其收敛意义}
\label{b2:subsec:B20203}

对 \(f\in L^2(-1,1)\)，有
\[f=\sum_{n=0}^\infty c_nP_n,\qquad
 c_n=\frac{2n+1}{2}\int_{-1}^1f(t)P_n(t)\,\mathrm dt,\]
且 Parseval 等式为 \(\lVert f\rVert_2^2=\sum 2|c_n|^2/(2n+1)\)。这是 \(L^2\) 收敛，不应自动解释为每点或端点的一致收敛。若具有更多光滑性，可以借微分方程把系数中的导数移到 \(f\) 上，取得更快衰减，但须同时核验边界项。

\ProseParagraphBreak

例如 \(f\in C^2[-1,1]\) 时，\((1-t^2)f'\) 在两端为零。对 \(n\ge1\)，分部积分给出
\[n(n+1)\int fP_n=\int Lf\,P_n.\]
因此归一化系数 \(\langle f,\sqrt{(2n+1)/2}P_n\rangle\) 不超过 \(\lVert Lf\rVert_2/[n(n+1)]\)。如果要反复使用这个等式，就应确认 \(L^jf\) 的端点行为和可积性，而不能只根据“\(f\) 有几阶导数”口头省去条件。

\ProseParagraphBreak

在三维球面上，轴对称函数写成 \(F(\theta)=f(\cos\theta)\)，面积积分为 \(2\pi\int_{-1}^1f(t)\,\mathrm dt\)。因此上述无权 Legendre 展开恰与轴对称球面展开匹配；若错误保留 \(\mathrm d\theta\) 而漏掉 \(\sin\theta\)，就会得到另一种正交问题。

% !TeX root = ../../Book2.tex
\section{球谐函数与特征函数展开}
\label{b2:sec:B203}

Legendre 多项式只描述三维球面上的轴对称部分。一般角向函数需要按所有方向分解；齐次调和多项式的球面限制提供了相应有限维模态。先从欧氏空间多项式构造这些空间，再证明正交性和完备性，就不必预先假定球面 Laplace 算子的谱已经齐全。

\subsection{极坐标下的 Laplace 算子}
\label{b2:subsec:B20301}

设 \(d\ge2\)，写 \(x=r\omega\)、\(r>0\)、\(\omega\in S^{d-1}\)。对球面光滑函数 \(Y\)，沿半径保持常值延拓到一个球壳，并把欧氏梯度在 \(r=1\) 的切向部分记为 \(\nabla_SY\)。球面散度与该切向梯度配对定义 \(\Delta_S\)。等价地，径向常值延拓 \(\widetilde Y\) 满足 \(\Delta_SY=(\Delta\widetilde Y)|_{r=1}\)。

\begin{proposition}[极坐标 Laplace 公式]
\label{b2:prop:B2-polar-laplacian}
对球壳内光滑函数 \(u(r,\omega)\)，
\[\Delta u=\partial_r^2u+\frac{d-1}r\partial_ru+\frac1{r^2}\Delta_Su.
\]
球面上有 Green 恒等式
\[\int_{S^{d-1}}(-\Delta_SY)\overline Z\,\mathrm d\sigma
 =\int_{S^{d-1}}\nabla_SY\cdot\overline{\nabla_SZ}\,\mathrm d\sigma.
\]
\end{proposition}
\begin{proof}
微分一个径向位移与切向位移，得到正交分解 \(\nabla u=\omega\partial_ru+r^{-1}\nabla_Su\)。在球壳上把 \(\int\nabla u\cdot\overline{\nabla v}\) 写为极坐标积分。径向项带权 \(r^{d-1}\)，分部积分产生 \(r^{1-d}\partial_r(r^{d-1}\partial_ru)\)；切向项带权 \(r^{d-3}\)，按球面散度定义分部积分产生 \(r^{-2}\Delta_Su\)。先对紧支撑于球壳的测试函数验证，得到分布等式；两边光滑，故逐点成立。对只依赖角向的函数重复同一计算，球面没有边界，切向分部积分即为所列 Green 恒等式。
\end{proof}

例如在 \(S^2\) 上，球坐标面积元为 \(\sin\theta\,\mathrm d\theta\,\mathrm d\varphi\)，因而
\[\Delta_S=\frac1{\sin\theta}\partial_\theta(\sin\theta\,\partial_\theta)
              +\frac1{\sin^2\theta}\partial_\varphi^2.
\]
若函数只依赖 \(t=\cos\theta\)，第一项化为 \((1-t^2)\partial_t^2-2t\partial_t\)，第二项为零，恢复 Legendre 方程中的微分表达式。坐标系数在极点发散，算子本身作为球面内在对象却没有奇点。

\subsection{齐次调和多项式}
\label{b2:subsec:B20302}

记 \(\mathcal P_\ell\) 为 \(d\) 个变量的 \(\ell\) 次齐次复多项式空间，\(\mathcal H_\ell=\{P\in\mathcal P_\ell:\Delta P=0\}\)。它们在单位球面的限制称为 \(\ell\) 阶 \term[qiuxie-hanshu]{球谐函数}[spherical harmonic]，相应空间记为 \(\mathcal Y_\ell\)。齐次性保证球面限制为零的多项式处处为零，所以 \(\mathcal H_\ell\) 与 \(\mathcal Y_\ell\) 同维。

\begin{proposition}[齐次多项式的调和分解]
\label{b2:prop:B2-harmonic-decomposition}
对 \(\ell\ge0\)，有唯一分解
\[\mathcal P_\ell=\mathcal H_\ell\oplus|x|^2\mathcal P_{\ell-2},\]
其中负阶空间取零。反复分解得到每个齐次多项式的唯一表达
\[P=\sum_{j=0}^{\lfloor\ell/2\rfloor}|x|^{2j}H_{\ell-2j},\qquad
 H_{\ell-2j}\in\mathcal H_{\ell-2j}.
\]
\end{proposition}
\begin{proof}
直接使用乘积法则及 Euler 齐次恒等式 \(x\cdot\nabla H_m=mH_m\)，对调和齐次 \(H_m\) 得到
\[\Delta(|x|^{2k}H_m)=2k(2m+d+2k-2)|x|^{2k-2}H_m.\]
对 \(\ell=0,1\) 分解显然。归纳假设 \(\mathcal P_{\ell-2}\) 已分解成 \(|x|^{2j}\mathcal H_{\ell-2-2j}\) 的直和。上式说明 \(\Delta\) 从 \(|x|^2\mathcal P_{\ell-2}\) 到 \(\mathcal P_{\ell-2}\) 在各分量上乘以严格正的常数，因此是线性同构。给定 \(P\in\mathcal P_\ell\)，唯一取 \(Q\in|x|^2\mathcal P_{\ell-2}\)，使 \(\Delta Q=\Delta P\)，便有 \(P-Q\in\mathcal H_\ell\)。同构的单射性又保证两个直和分量交集为零，完成归纳。
\end{proof}

于是
\[N_\ell=\dim\mathcal Y_\ell
 =\binom{\ell+d-1}{d-1}-\binom{\ell+d-3}{d-1},\]
第二项在 \(\ell<2\) 时取零。三维为 \(N_\ell=2\ell+1\)，二维除 \(N_0=1\) 外均为 \(2\)，对应圆周上的正弦、余弦两种模态。低阶例子中，常数是零阶，坐标函数是一阶；二阶调和多项式是迹为零的二次型，如 \(x_1^2-x_2^2\) 和 \(x_1x_2\)。

\subsection{球面谱与完备性}
\label{b2:subsec:B20303}

若 \(P(r\omega)=r^\ell Y(\omega)\) 且 \(P\) 调和，极坐标公式给出
\[-\Delta_SY=\ell(\ell+d-2)Y.\]
不同 \(\ell\) 的特征值严格递增，因此由球面 Green 恒等式，\(\mathcal Y_\ell\) 两两正交。每个空间选择一组标准正交基 \(Y_{\ell,1},\ldots,Y_{\ell,N_\ell}\)。

\begin{theorem}[球谐展开]
\label{b2:thm:B2-spherical-expansion}
上述全部球谐基组成 \(L^2(S^{d-1},\mathrm d\sigma)\) 的标准正交基。对 \(f\in L^2(S^{d-1})\)，
\[f=\sum_{\ell=0}^\infty\sum_{j=1}^{N_\ell}c_{\ell,j}Y_{\ell,j},\qquad
 c_{\ell,j}=\int_{S^{d-1}}f\overline{Y_{\ell,j}}\,\mathrm d\sigma,
\]
且 \(\lVert f\rVert_2^2=\sum_{\ell,j}|c_{\ell,j}|^2\)，级数按 \(L^2\) 意义收敛。
\end{theorem}
\begin{proof}
正交性已证，只须证明有限球谐和稠密。任一多项式可拆成齐次部分，每个齐次部分再按调和分解展开；在球面上 \(|x|=1\)，所以其限制是有限球谐和。

多项式限制在球面连续函数中一致稠密。为给出具体依据，先把连续球面函数沿半径延拓到球壳，再乘一个在原点附近为零的径向截断，得到包住球面的立方体上的连续函数。将立方体缩放到 \([0,1]^d\)，采用各坐标 Bernstein 权重的乘积。权重和为 \(1\)，各坐标的平均平方偏差至多为 \(1/(4n)\)，故偏离给定点超过 \(\varepsilon\) 的总权重至多为 \(d/(4n\varepsilon^2)\)。结合一致连续性，逼近误差至多为连续模加上 \(2\lVert f\rVert_\infty d/(4n\varepsilon^2)\)，先取小 \(\varepsilon\) 再取大 \(n\)，即得一致逼近。限制回球面便得到所需多项式逼近。

球面面积是有限 Radon 测度，连续函数在其 \(L^2\) 中稠密。因此有限球谐和在 \(L^2\) 中稠密，标准正交展开定理给出展开与 Parseval 等式。
\end{proof}

这里完备性由多项式分解与逼近证明，并未从已经假定完整的球面谱反推。对光滑 \(f\)，还可多次用 Green 恒等式，得到
\[\langle(-\Delta_S)^kf,Y_{\ell,j}\rangle
 =[\ell(\ell+d-2)]^kc_{\ell,j}.\]
Parseval 因而给出加权系数平方可和。这样的估计描述高阶角向光滑性；普通 \(L^2\) 函数则只具有未加权的平方可和，不能据此随意逐项求导。

\subsection{加法公式与球内调和延拓}
\label{b2:subsec:B20304}

三维时，正交投影到 \(\mathcal Y_\ell\) 的核为
\[Z_\ell(\omega,\eta)=\sum_{j=1}^{2\ell+1}Y_{\ell,j}(\omega)\overline{Y_{\ell,j}(\eta)}.
\]
核不依赖所选正交基，并同时旋转不变。固定 \(\eta\) 后，绕该轴的旋转不变性使它只依赖 \(t=\omega\cdot\eta\)；它还是 \(\ell\) 阶球谐多项式，故满足 Legendre 方程。规则多项式解是一维的，由最高次系数递推可见，所以它是 \(P_\ell(t)\) 的常数倍。对角值由旋转不变性为常数，且积分等于投影的维数 \(2\ell+1\)，得到 \term[qiuxie-jiafa-gongshi]{球谐加法公式}[spherical harmonic addition formula]
\[Z_\ell(\omega,\eta)=\frac{2\ell+1}{4\pi}P_\ell(\omega\cdot\eta).
\]
这一推导同时解释了 Legendre 多项式为什么只需一个变量，却能汇集整个 \(2\ell+1\) 维角向空间。

\ProseParagraphBreak

在一般维数，对角核同理为 \(\sum_j|Y_{\ell,j}(\omega)|^2=N_\ell/|S^{d-1}|\)。因此给定 \(f\in L^2(S^{d-1})\)，级数
\[u(r\omega)=\sum_{\ell,j}r^\ell c_{\ell,j}Y_{\ell,j}(\omega),\qquad0\le r<1,\]
在每个 \(r\le\rho<1\) 上一致收敛。事实上，Cauchy--Schwarz 不等式给出
\[
 \left|\sum_jr^\ell c_{\ell,j}Y_{\ell,j}(\omega)\right|
 \le r^\ell\left(\sum_j|c_{\ell,j}|^2\right)^{1/2}
       \sqrt{\frac{N_\ell}{|S^{d-1}|}}.
\]
维数只按多项式增长，而 \(\rho^\ell\) 几何衰减，因此各阶项的上界可求和。各部分和是调和多项式，一致局部极限保留均值性质，故 \(u\) 在球内调和。

\ProseParagraphBreak

又由 Parseval，\(\lVert u(r\cdot)-f\rVert_{L^2(S^{d-1})}^2=\sum_{\ell,j}|r^\ell-1|^2|c_{\ell,j}|^2\to0\)，支配收敛保证边界数据按 \(L^2\) 意义恢复。这个构造与第十三章球内 Poisson 核相容：两者对光滑边界数据给出同一调和函数，再由 \(L^2\) 连续性延拓。若要求逐点边界收敛，还需额外的边界点或数据假设，不能只凭这个平方范数极限得到。

\ProseParagraphBreak

对于 Helmholtz 方程，代入单个角向模态 \(u(r,\omega)=v(r)Y_{\ell,j}(\omega)\)，径向方程为
\[v''+\frac{d-1}rv'+\left(\kappa^2-\frac{\ell(\ell+d-2)}{r^2}\right)v=0.
\]
令 \(v=r^{-(d-2)/2}w(\kappa r)\)，便得到阶数 \(\ell+(d-2)/2\) 的 Bessel 方程。于是本附录的角向展开与附录 A 的径向函数共同完成球形几何中的分离变量，阶数中多出的 \((d-2)/2\) 来自径向体积因子，而不是任意选择。


\begin{chaptersummary}
Hermite 体系把多项式的代数与 Gauss 衰减结合，提供谐振子的离散谱及 Fourier 变换的特征基。Legendre 体系通过端点消失的 Rodrigues 因子获得无权正交性，再由对称微分算子推出方程与递推。两种体系的完备性都另有证明，不能从形式特征关系自动得到。

\ProseParagraphBreak

球谐空间来自齐次多项式的唯一调和分解，其正交性由球面 Green 公式给出，完备性由多项式逼近给出。角向特征值与径向体积项共同确定 Bessel 阶数；加法公式则把一个角向特征空间压缩为旋转不变的投影核。所有展开首先具有明确的 \(L^2\) 意义，更强收敛及逐项求导须另用正则性估计。
\end{chaptersummary}
\BookExercises

各题均为本书自拟，要求把本附录的基函数与投影、数据和具体算子结合。

% 自拟：Fourier四个特征空间的投影。
\begin{exercise}\label{b2:ex:B2-fourier-projections}
在 \(L^2(\mathbb R)\) 上，对 \(k=0,1,2,3\) 定义 \(P_k=\tfrac14\sum_{j=0}^3 i^{kj}\mathcal F^j\)。证明它是到 \(\operatorname{span}\{\psi_n:n\equiv k\pmod4\}\) 的闭包上的正交投影，并写出 \(\sum_kP_k\) 与 \(P_kP_l\)。
\end{exercise}
\begin{solution}
对 \(\psi_n\)，有 \(P_k\psi_n=\tfrac14\sum_{j=0}^3i^{(k-n)j}\psi_n\)，系数在 \(n\equiv k\pmod4\) 时为 \(1\)，否则为零。由于 \(\psi_n\) 是标准正交基且各算子有界，结论推广到全部 \(L^2\)。它在系数上保留对应剩余类，故自伴且幂等，\(\sum_kP_k=I\)、\(P_kP_l=0\)（\(k\ne l\)）。
\end{solution}

% 自拟：谐振子的基态与不确定性下界。
\begin{exercise}\label{b2:ex:B2-ground-state}
设 \(f\in\mathcal S(\mathbb R)\)。证明
\[\int_\mathbb R(|f'|^2+x^2|f|^2)\,\mathrm dx\ge\lVert f\rVert_2^2,\]
并刻画等号。再求 \(-\alpha\frac{\mathrm d^2}{\mathrm dx^2}+\beta x^2\)、\(\alpha,\beta>0\) 的最小特征值与归一化基态。
\end{exercise}
\begin{solution}
因 \(\mathcal L=2a^*a+I\)，左端等于 \(2\lVert af\rVert_2^2+\lVert f\rVert_2^2\)。等号当且仅当 \(f'+xf=0\)，即 \(f=Ce^{-x^2/2}\)。对一般系数令 \(s=(\beta/\alpha)^{1/4}\)、\(y=sx\)，算子变成 \(\sqrt{\alpha\beta}(-D_y^2+y^2)\)。最小特征值为 \(\sqrt{\alpha\beta}\)，归一化基态是 \(s^{1/2}\pi^{-1/4}e^{-s^2x^2/2}\)，相差模为 \(1\) 的常数不影响归一化。
\end{solution}

% 自拟：热半群的Hermite系数构造与误差。
\begin{exercise}\label{b2:ex:B2-oscillator-semigroup}
对 \(f=\sum c_n\widetilde h_n\in L^2\) 定义 \(T_tf=\sum e^{-(2n+1)t}c_n\widetilde h_n\)、\(t\ge0\)。证明半群关系、强连续性以及 \(\lVert T_t\rVert=e^{-t}\)，并证明当 \(t>0\) 时像属于 \(D(\mathcal L^m)\) 对每个整数 \(m\ge1\) 成立。
\end{exercise}
\begin{solution}
系数相乘给出 \(T_sT_t=T_{s+t}\)。由 \(|e^{-(2n+1)t}-1|\le2\) 及 \(\sum|c_n|^2<\infty\)，支配收敛给出 \(T_tf\to f\)。算子范数是系数模的上确界 \(e^{-t}\)，并在 \(\widetilde h_0\) 上取等。固定 \(t>0\) 时，\(\sup_{\lambda\ge1}\lambda^{2m}e^{-2t\lambda}<\infty\)，故加权系数平方和有限，得到所列定义域性质。这是谱意义的平滑作用。
\end{solution}

% 自拟：有限Legendre投影的显式误差。
\begin{exercise}\label{b2:ex:B2-legendre-projection}
把 \(t^4\) 展成 \(P_0,P_2,P_4\) 的线性组合，求它在次数至多为 \(2\) 的多项式空间中的最佳 \(L^2(-1,1)\) 逼近及误差范数。
\end{exercise}
\begin{solution}
由 \(P_2=(3t^2-1)/2\)、\(P_4=(35t^4-30t^2+3)/8\)，解得
\[t^4=\frac15P_0+\frac47P_2+\frac8{35}P_4.\]
前两项是到二次多项式空间的正交投影（奇项系数为零）。误差范数为 \(\frac8{35}\sqrt{2/9}=\frac{8\sqrt2}{105}\)。最优性来自正交投影定理，而不只是比较某几个候选多项式。
\end{solution}

% 自拟：三维二次与四次齐次多项式的调和分解。
\begin{exercise}\label{b2:ex:B2-harmonic-parts}
在 \(\mathbb R^3\) 中分别分解 \(x_3^2\) 与 \(x_3^4\) 为 \(\sum_j|x|^{2j}H_{\ell-2j}\)，并说明限制到球面后与 Legendre 展开的关系。
\end{exercise}
\begin{solution}
记 \(r=|x|\)。二次分解为 \(x_3^2=(x_3^2-r^2/3)+r^2/3\)。四次分解为
\[x_3^4=\left(x_3^4-\frac67r^2x_3^2+\frac3{35}r^4\right)
 +\frac67r^2\left(x_3^2-\frac13r^2\right)+\frac15r^4.
\]
逐项计算 Laplace，或使用正文的 \(\Delta(r^{2k}H_m)\) 公式，验证括号内分别为四阶、二阶调和多项式。限制到 \(r=1\) 后，它们分别与 \(P_4(x_3)\)、\(P_2(x_3)\) 成比例，恰恢复上一题的 Legendre 系数。
\end{solution}

% 自拟：球面均值、一次模态及最佳近似。
\begin{exercise}\label{b2:ex:B2-sphere-projection}
对 \(f\in L^2(S^{d-1})\)，求它在常数与坐标函数张成的空间中的正交投影。写出系数的积分公式，并解释余项的零均值和零一阶矩条件。
\end{exercise}
\begin{solution}
由对称性，常数与各坐标正交，\(\int\omega_i\omega_j=(|S^{d-1}|/d)\delta_{ij}\)。因此投影为 \(a+b\cdot\omega\)，其中
\[a=\frac1{|S^{d-1}|}\int f,\qquad b_i=\frac d{|S^{d-1}|}\int f(\omega)\omega_i\,\mathrm d\sigma.\]
余项与零阶、一次球谐空间均正交，恰好表示其均值及全部一阶矩为零。对复值 \(f\)，同一公式成立，因为坐标函数为实值，内积仍对第一个变量线性。
\end{solution}

% 自拟：利用球谐次数给出非平凡边界多项式延拓。
\begin{exercise}\label{b2:ex:B2-ball-quartic}
在三维单位球上求边界数据 \(f(\omega)=\omega_3^4\) 的调和延拓，写成直角坐标多项式，并求球心值。
\end{exercise}
\begin{solution}
上面的调和分解在球面给出四阶、二阶和常数三项。向球内延拓时分别保留对应齐次调和多项式，得到
\[u(x)=x_3^4-\frac67|x|^2x_3^2+\frac3{35}|x|^4
       +\frac67\left(x_3^2-\frac13|x|^2\right)+\frac15.
\]
它调和且边界值正确，由球内连续 Dirichlet 唯一性即为所求。中心值为 \(1/5\)，也等于球面上 \(\omega_3^4\) 的归一化平均。
\end{solution}

% 自拟：角向模态改变径向Bessel阶数与边界频率。
\begin{exercise}\label{b2:ex:B2-helmholtz-ball}
对单位球内 \(-\Delta u=\kappa^2u\)、零边界条件，固定 \(\ell\) 阶球谐函数 \(Y\)。求原点正则的分离变量解形式，并说明允许的正 \(\kappa\) 满足什么方程；不要据此直接宣称整个 Dirichlet 谱已经完备。
\end{exercise}
\begin{solution}
径向化给出阶数 \(\mu=\ell+(d-2)/2\) 的 Bessel 方程。原点正则支为
\[u(r,\omega)=C r^{-(d-2)/2}J_\mu(\kappa r)Y(\omega),\]
其领先行为为常数乘 \(r^\ell Y\)，与齐次调和多项式相容；第二类支在原点具有不允许的奇性。零边界要求 \(J_\mu(\kappa)=0\)。这构造了一族特征函数，但径向不同零点模态的完备性仍需奇异 Sturm--Liouville 或球上自伴紧逆理论，不能只由球谐角向完备性推出。
\end{solution}
